\documentclass[a4paper,12pt]{ltjsarticle}
\usepackage[margin=25mm]{geometry}
\usepackage{amsmath,amssymb}
\usepackage{enumitem}

\begin{document}

\begin{center}
{\LARGE 数I・A 共通テスト風（選択肢なし・大問形式）}
\end{center}

\vspace{1em}

\section*{第1問（数と式・二次関数）}
実数 $x$ について、関数
\[
f(x)=x^2-6x+a \quad (a\text{は定数})
\]
を考える。
\begin{enumerate}[label=\arabic*.]
  \item $f(x)$ の最小値を $a$ を用いて表せ。
  \item 方程式 $x^2-6x+a=0$ が異なる2つの実数解をもつような $a$ の範囲を求めよ。
  \item さらに、2つの解を $\alpha,\beta$（$\alpha<\beta$）とするとき、$\alpha<1<\beta$ が成り立つような $a$ の範囲を求めよ。
  \item $a=5$ のとき、$f(x)\le 0$ を満たす $x$ の範囲を求めよ。
\end{enumerate}

\section*{第2問（場合の数・確率）}
1から8までの整数が1枚ずつ書かれた8枚のカードがある。この中から同時に3枚を取り出す試行を考える。
\begin{enumerate}[label=\arabic*.]
  \item 3枚の和が偶数となる確率を求めよ。
  \item 3枚のうち最大の数が6である確率を求めよ。
  \item 3枚の積が3の倍数となる確率を求めよ。
  \item 3枚を小さい順に $a<b<c$ とするとき、$a+b>c$ を満たす確率を求めよ。
\end{enumerate}

\section*{第3問（図形と計量）}
三角形 $ABC$ において、
\[
AB=6,\quad AC=8,\quad \angle A=60^\circ
\]
とする。
\begin{enumerate}[label=\arabic*.]
  \item 三角形 $ABC$ の面積を求めよ。
  \item 外接円の半径を $R$ とするとき、$R$ を求めよ。
  \item 辺 $BC$ を底辺としたときの高さを求めよ。
\end{enumerate}

\section*{第4問（データの分析）}
あるクラス40人の数学テスト（100点満点）の結果について、平均点が64点、標準偏差が12点であった。  
このとき、得点を $x$ 点とする生徒の偏差値 $H$ を
\[
H=50+10\frac{x-64}{12}
\]
で定める。
\begin{enumerate}[label=\arabic*.]
  \item $x=76$ のときの偏差値を求めよ。
  \item 偏差値が60以上であるための $x$ の範囲を求めよ。
\end{enumerate}

\end{document}
